% =============================================================================
% chapters/ch1-introduction.tex: Chapter 1.
%
% HOW A CHAPTER FILE WORKS (applies to every file in chapters/)
% - The file is pulled in by \include in main.tex, which starts a new page.
% - \chapter{Title} prints "Chapter 1" + title and starts the running header; \label{ch:...} right after it lets you write \cref{ch:introduction} ("Chapter 1") anywhere.
% - \section / \subsection are numbered (1.1, 1.1.1) and listed in the table of contents; \subsubsection is unnumbered. Add \label{sec:...} after each one you want to refer to.
% - A heading that needs an abbreviation: \section{Why \glsfmtshort{mcp} ...} (plain "MCP" also works but won't count as a use of the glossary entry).
% - Starred versions (\section*{...}) are unnumbered and not in the contents.
% - Citations: \cite{kumar2018rev2}, \cite[p.~3]{key} (keys: references.bib).
%   Glossary: \gls{mcp} (entries: glossary/entries.tex).
%   Figures, tables, algorithms: see README.md.
% - A blank line starts a new paragraph; a single line break does not.
% - % starts a comment; to print a percent sign write \%.
%   Other characters that need a backslash: \& \# \_ \$ \{ \}.
% =============================================================================

\chapter{Introduction}\label{ch:introduction}
